% ============================================================
% paper/sections/12_conclusion_scope_boundary.tex
% Conclusion and Scope Boundary
% ============================================================

\section{Conclusion and Scope Boundary}

\subsection{Summary of the Contribution}

This whitepaper has presented a reviewer-grade, executable-closed,
diagnostic-only specification of the Enterprise Continuum
$K_{\mathrm{EA}}$, projected directly from the immutable Executable
Theory Specification \texttt{ETS.K\_EA.v1.1}.

The document has:
\begin{itemize}
  \item positioned enterprise diagnosis within the Ontology-of-Continua
        framework,
  \item defined the enterprise as a conditional, bounded continuum,
  \item formalized admissibility, persistence, and termination
        conditions,
  \item treated STOP as a first-class diagnostic outcome,
  \item preserved invariants and No-Go results as structural guarantees.
\end{itemize}

No new primitives, parameters, or claims have been introduced.

\subsection{What This Whitepaper Enables}

Within its scope, this whitepaper enables:
\begin{itemize}
  \item structurally consistent enterprise diagnosis,
  \item explicit identification of non-diagnosable regimes,
  \item rigorous distinction between degradation and structural failure,
  \item inspection and review of diagnostic assumptions.
\end{itemize}

These capabilities arise solely from the formal structure of the
continuum.

\subsection{What This Whitepaper Explicitly Forbids}

The document explicitly forbids:
\begin{itemize}
  \item prescriptive interpretation of diagnostic results,
  \item optimization or improvement guidance,
  \item substitution of structural diagnosis with isolated metrics,
  \item bypassing STOP, invariants, or No-Go results.
\end{itemize}

Any such use would exceed the admissible scope of the model.

\subsection{Boundary of Applicability}

The applicability of $K_{\mathrm{EA}}$ is bounded by the conditions under
which the enterprise can be treated as a diagnosable continuum.

When existence, observability, persistence, or parent compatibility
conditions are violated, diagnosis terminates. The model makes no claims
beyond this boundary.

\subsection{Final Remark}

The Enterprise Continuum is not a promise of success or recovery. It is
a formal object designed to preserve diagnostic integrity by making its
own limits explicit.

Within those limits, it provides a coherent and reproducible basis for
enterprise diagnosis. Outside them, it deliberately remains silent.
